% Project 01 report. Build from this folder:  latexmk -pdf main.tex
% Limit: 6-8 pages, double column. Target ~7.5 pages. Page budget per section is
% noted in each section header comment.
% Red [TODO: ...] marks are placeholders; none may remain in the submitted PDF.
\documentclass[conference]{IEEEtran}
\usepackage[T1]{fontenc}
\usepackage{amsmath}
\usepackage{newtxtext,newtxmath}   % Times for text and math, as in IEEE papers
\usepackage{graphicx}
\usepackage{booktabs}
\usepackage{xcolor}
\usepackage{cite}
\usepackage{url}
\usepackage[hidelinks]{hyperref}
\graphicspath{{../results/}}

\newcommand{\todo}[1]{\textcolor{red}{[TODO: #1]}}
% Result figures without their matplotlib title (the caption replaces it).
\newcommand{\resfig}[1]{\includegraphics[width=\textwidth,trim=0 0 0 20bp,clip]{#1}}

\begin{document}

\title{Tabular Q-Learning versus Deep Q-Networks on Gymnasium Control Tasks:
Discretization, Hyperparameter Sensitivity, and Ablations}

\author{
\IEEEauthorblockN{Jos\'e Reinoso}
\IEEEauthorblockA{School of Mathematical and Computational Sciences\\
Yachay Tech University, Urcuqu\'i, Ecuador}
\and
\IEEEauthorblockN{Aldrin Chavez}
\IEEEauthorblockA{School of Mathematical and Computational Sciences\\
Yachay Tech University, Urcuqu\'i, Ecuador}
}

\maketitle

% ---------------------------------------------------------------- 0.3 page
\begin{abstract}
We compare tabular reinforcement learning with Deep Q-Networks (DQN) on Gymnasium
environments under a common protocol: equal budgets of environment steps, several
independent seeds, and greedy evaluation reported separately from training returns. On
CartPole-v1, tabular Q-learning with uniform state discretization (6, 10 and 20 bins per
dimension and three velocity ranges) and SARSA are compared with DQN under the same
200k-step budget and five seeds. On LunarLander-v3 we study the sensitivity of DQN to the
learning rate and to the target-network update period, ablate the target network and the
replay buffer, and compare Double and Dueling DQN with vanilla DQN, with three seeds per
configuration. On CartPole, discretized Q-learning reaches a return of about 200 within 20k
steps but never solves the task, whereas DQN reaches the 475 threshold in all five seeds
(median 90k steps) at 18 times the wall-clock cost per step, and then degrades to a final
return of 159. On LunarLander, the default DQN ends at $207\pm30$; removing the replay buffer
($-76$) hurts more than removing the target network (83), and the largest learning rate
makes one seed diverge. Double and Dueling DQN (214--223) stay within the across-seed spread
of vanilla DQN. A written analysis quantifies, through action- and state-space counts, why continuous-action
locomotion tasks and Atari pixel observations are out of reach for tabular methods and DQN,
and which algorithm families suit them.
\end{abstract}

\begin{IEEEkeywords}
reinforcement learning, Q-learning, SARSA, deep Q-network, discretization, ablation study
\end{IEEEkeywords}

% ---------------------------------------------------------------- 0.5 page
\section{Introduction}

Value-based reinforcement learning (RL) estimates the action-value function $Q(s,a)$ and
acts greedily with respect to it. In its tabular form, Q-learning~\cite{watkins1992} stores
one value per state--action pair, which requires a finite and small state
space~\cite{sutton2018}. Most control problems, however, have continuous observations.
Two routes are then available: discretize the state space so that a table can still be
used, or replace the table with a function approximator, as Deep Q-Networks (DQN) do with a
neural network, an experience-replay buffer and a target network~\cite{mnih2015}.
Discretization is simple and cheap per step, but the number of states grows exponentially
with the number of dimensions~\cite{bellman1957}. DQN generalizes across states, but
bootstrapping with a nonlinear approximator trained off-policy can be
unstable~\cite{sutton2018}.

This paper compares both routes on Gymnasium~\cite{towers2024} tasks under a common
protocol that follows the recommendations of~\cite{henderson2018}: equal budgets of
environment steps, multiple seeds, and greedy evaluation separated from training returns.
Our contributions are:
\begin{itemize}
  \item a comparison of tabular Q-learning (three bin counts, three velocity ranges, and
        SARSA) with DQN on CartPole-v1 under the same 200k-step budget (Tier~1);
  \item a one-at-a-time study of the learning rate and the target-update period of DQN on
        LunarLander-v3, with ablations of the target network and of the replay buffer
        (Tier~2);
  \item a comparison of Double~\cite{vanhasselt2016} and Dueling~\cite{wang2016} DQN with
        vanilla DQN on the same task (Tier~3, option~B);
  \item a quantitative analysis of why continuous-action locomotion and pixel-based Atari
        games fall outside the reach of tabular methods and DQN.
\end{itemize}
Section~\ref{sec:pertinence} justifies the choice of environments,
Sections~\ref{sec:methods} and~\ref{sec:setup} describe the methods and the protocol,
Sections~\ref{sec:results} and~\ref{sec:discussion} report and discuss the results, and
Section~\ref{sec:analysis} contains the analysis of unsuitable environments. Code,
configurations and logs are available at
\url{https://github.com/gato09sa/rl-dqn-project01}.

% ---------------------------------------------------------------- 0.4 page
\section{Pertinence Analysis}\label{sec:pertinence}

A task suits tabular RL if its state space is finite and small, or can be discretized into
a table whose entries are visited often enough within the training budget. It suits DQN if
its action space is discrete, because the greedy action is an $\arg\max$ over one network
output per action; the observation space may be continuous. Table~\ref{tab:envs} applies
these criteria to the candidate tasks.

\begin{table}[t]
\centering
\caption{Candidate environments (Gymnasium 1.3.0). ``10-bin $|\mathcal{S}|$'': number of
table states with 10 bins per continuous dimension.}
\label{tab:envs}
\footnotesize
\setlength{\tabcolsep}{4pt}
\begin{tabular}{@{}lcccrl@{}}
\toprule
Environment & Obs. & $|\mathcal{A}|$ & Threshold & 10-bin $|\mathcal{S}|$ & Suits \\
\midrule
CartPole-v1    & 4             & 2 & 475    & $10^{4}$           & tabular, DQN \\
MountainCar-v0 & 2             & 3 & $-110$ & $10^{2}$           & tabular, DQN \\
Acrobot-v1     & 6             & 3 & $-100$ & $10^{6}$           & DQN \\
LunarLander-v3 & 8$^{\dagger}$ & 4 & 200    & $4\!\times\!10^{6}$ & DQN \\
\bottomrule
\multicolumn{6}{@{}l}{\scriptsize $^{\dagger}$Six continuous dimensions and two binary
leg-contact flags.}
\end{tabular}
\end{table}

\textbf{CartPole-v1} has four continuous observation dimensions, two of which (cart and pole
velocities) are unbounded, so bin ranges must be chosen by hand. With 10 bins per
dimension the table has $10^4$ states and $2\times10^4$ entries, ten times fewer than the
$2\times10^5$ training steps, so tabular learning is feasible. With 20 bins it has
$1.6\times10^5$ states and $3.2\times10^5$ entries, more than the step budget, so many
entries can be updated only rarely, if at all. The two actions are discrete, so DQN applies
directly. CartPole therefore suits both families and is our Tier~1 task.

\textbf{LunarLander-v3} has six continuous dimensions, two binary flags and four discrete
actions. Ten bins per continuous dimension give $4\times10^6$ states and $1.6\times10^7$
entries, 40 times the 400k-step budget: the table cannot be filled, whereas DQN only needs
the discrete action space. It is therefore a DQN task, used for Tiers~2 and~3.
MountainCar-v0 (Tier~1 alternative) and Acrobot-v1 (Tier~2 alternative) were not used.
MountainCar's sparse reward makes a naive DQN rarely reach the goal under
$\varepsilon$-greedy exploration, so the comparison would mostly measure exploration
rather than the effect of discretization, which is the question Tier~1 asks; on CartPole
both families can reach the threshold, which makes their comparison informative.
Acrobot is the option offered for weak hardware; a timing test on one of our laptops
(Intel i5-10210U, four runs in parallel) showed LunarLander fits the CPU budget, and its
richer dynamics leave more room for the hyperparameters, the ablations and the Tier~3
variants to differ.

% ---------------------------------------------------------------- 1.0 page
\section{Methods}\label{sec:methods}

\subsection{Tabular Q-learning and SARSA}
Both methods update a table $Q(s,a)$ after every transition $(s_t,a_t,r_{t+1},s_{t+1})$:
\begin{equation}
Q(s_t,a_t) \leftarrow Q(s_t,a_t) + \alpha\,\bigl[y_t - Q(s_t,a_t)\bigr],
\end{equation}
with target $y_t = r_{t+1} + \gamma \max_{a} Q(s_{t+1},a)$ for Q-learning
(off-policy)~\cite{watkins1992} and $y_t = r_{t+1} + \gamma\,Q(s_{t+1},a_{t+1})$ for SARSA
(on-policy), where $a_{t+1}$ is the action the agent actually takes
next~\cite{rummery1994}. If $s_{t+1}$ is terminal, $y_t = r_{t+1}$. An episode cut by the
time limit is \emph{not} terminal, so its last transition still bootstraps from
$s_{t+1}$~\cite{pardo2018}. The table is initialized to zero, actions are chosen
$\varepsilon$-greedily, and ties in the greedy action are broken at random so that
unvisited states do not bias the policy toward action~0.

\subsection{State discretization}
Each observation dimension $i$ is split into $b_i$ equal-width bins over a range
$[l_i,h_i]$; values outside the range fall into the first or last bin. The per-dimension
indices $\phi_i \in \{0,\dots,b_i-1\}$ are combined into a single state index
$s = \sum_i \phi_i \prod_{j>i} b_j$, so $|\mathcal{S}| = \prod_i b_i$. For CartPole, the
cart position and pole angle use their termination limits ($\pm2.4$ and $\pm0.21$\,rad);
the unbounded velocities use $\pm3.0$ and $\pm3.5$ by default, $\pm1.5$ and $\pm2.0$ in
the \emph{narrow} variant, and $\pm6.0$ and $\pm7.0$ in the \emph{wide} variant.

\subsection{Deep Q-Network}
A multilayer perceptron with ReLU activations maps an observation to one value per action,
$Q(s,\cdot\,;\theta)$. Transitions are stored in a FIFO replay buffer $\mathcal{D}$ of
capacity $N$; after $T_0$ warm-up steps, one gradient step is taken every $k$ environment
steps on a minibatch of size $B$ sampled uniformly from $\mathcal{D}$, minimizing
\begin{align}
\mathcal{L}(\theta) &= \mathbb{E}_{\mathcal{D}}\Bigl[\ell_{H}\bigl(y - Q(s,a;\theta)\bigr)\Bigr], \\
y &= r + \gamma\,(1-d)\max_{a'} Q(s',a';\theta^{-}), \label{eq:dqn}
\end{align}
where $\ell_H$ is the Huber loss~\cite{huber1964}, $d$ is 1 only for true terminal states,
and $\theta^-$ are the parameters of a target network copied from $\theta$ every $C$
environment steps. We use Adam~\cite{kingma2015} with the gradient norm clipped at 10. The
exploration rate $\varepsilon$ decays linearly from 1.0 to 0.05 over $T_\varepsilon$ steps
and then stays constant. All networks are implemented in PyTorch~\cite{paszke2019}.

\subsection{Ablations}
\emph{No target network}: $\theta^- \equiv \theta$, i.e., the bootstrap target in
\eqref{eq:dqn} is computed by the online network itself (without gradient).
\emph{No replay buffer}: the buffer holds only the last $B$ transitions and each update uses
exactly this window, so minibatches are consecutive, correlated samples; the number of
gradient steps is the same as in vanilla DQN.

\subsection{Double and Dueling DQN (Tier~3)}
Double DQN~\cite{vanhasselt2016} decouples action selection from evaluation to reduce the
overestimation caused by the $\max$ operator:
$y = r + \gamma(1-d)\,Q\bigl(s',\arg\max_{a'}Q(s',a';\theta);\theta^-\bigr)$.
Dueling DQN~\cite{wang2016} splits the last layer into a state value and action advantages,
$Q(s,a) = V(s) + A(s,a) - \frac{1}{|\mathcal{A}|}\sum_{a'}A(s,a')$; subtracting the mean
advantage makes $V$ and $A$ identifiable.

% ---------------------------------------------------------------- 0.8 page
\section{Experimental Setup}\label{sec:setup}

\textbf{Experiments.} Tier~1 (CartPole-v1, 200k steps, 5 seeds) trains seven variants:
Q-learning with 6, 10 and 20 bins, Q-learning with 10 bins and narrow or wide velocity
ranges, SARSA with 10 bins, and DQN. Tier~2 (LunarLander-v3, 400k steps, 3 seeds) trains
the default DQN, two alternative learning rates ($10^{-4}$, $2\times10^{-3}$), two
alternative target-update periods ($C = 250$, $4000$), and the two ablations, changing one
value at a time. Tier~3 trains Double, Dueling and Double+Dueling DQN with the Tier~2
defaults and seeds, and reuses the Tier~2 default runs as the vanilla baseline.
Table~\ref{tab:hparams} lists all hyperparameters.

\begin{table}[t]
\centering
\caption{Hyperparameters of the default configurations.}
\label{tab:hparams}
\footnotesize
\setlength{\tabcolsep}{3pt}
\begin{tabular}{@{}lccc@{}}
\toprule
 & \multicolumn{2}{c}{Tier 1 (CartPole)} & Tiers 2--3 \\
\cmidrule(lr){2-3}
Parameter & Tabular & DQN & DQN (LunarLander) \\
\midrule
Step budget / seeds          & 200k / 5     & 200k / 5     & 400k / 3 \\
Evaluation period            & 5k           & 5k           & 10k \\
Discount $\gamma$            & 0.99         & 0.99         & 0.99 \\
$\varepsilon$: start $\to$ end & $1\to0.05$ & $1\to0.05$   & $1\to0.05$ \\
Decay steps $T_\varepsilon$  & 100k         & 20k          & 100k \\
Step size $\alpha$           & 0.1          & --           & -- \\
Bins per dimension           & 10 (6, 20)   & --           & -- \\
Hidden layers                & --           & 128, 128     & 256, 256 \\
Learning rate                & --           & $5\!\times\!10^{-4}$ & $5\!\times\!10^{-4}$\,$^{a}$ \\
Batch size $B$               & --           & 64           & 128 \\
Replay capacity $N$          & --           & 50k          & 100k \\
Warm-up steps $T_0$          & --           & 1k           & 10k \\
Update every $k$ steps       & 1            & 1            & 4 \\
Target period $C$ (steps)    & --           & 1000         & 1000\,$^{b}$ \\
\bottomrule
\multicolumn{4}{@{}l}{\scriptsize Also studied: $^{a}\,10^{-4}$ and $2\times10^{-3}$;
$^{b}\,250$ and $4000$ steps.}
\end{tabular}
\end{table}

\textbf{Evaluation protocol.} Every 5k (Tier~1) or 10k (Tiers~2--3) environment steps the
current greedy policy ($\varepsilon=0$) is evaluated on 10 episodes with fixed environment
seeds; training returns are logged separately, one value per episode. After training, the
final greedy policy is evaluated on 100 episodes with unseen seeds, which matches the
definition of the success thresholds (mean return over 100 episodes). \emph{Steps to
threshold} is the first periodic evaluation whose 10-episode mean reaches the threshold.
All curves use environment steps on the $x$-axis, so tabular agents and DQN are compared
at equal experience. They show the mean over seeds with a band of $\pm1$ population standard
deviation; training curves are smoothed with a 10-episode moving average and interpolated
onto a common step grid.

\textbf{Exploration schedules.} In Tier~1 the tabular agents decay $\varepsilon$ over half
the budget and DQN over 20k steps. A table has no generalization, so every state has to be
visited; DQN generalizes and benefits from exploiting earlier. The step budget, which is
the quantity compared, is identical.

\textbf{Reproducibility and hardware.} Python, NumPy, PyTorch and the environments are
seeded per run. Runs execute as four parallel processes with one CPU thread each; every run
stores its configuration, hardware and library versions. Tier~1 ran on an Intel Core
i5-10210U with Python 3.14.4 and PyTorch 2.14.1, and Tiers~2--3 on an Intel Core i5-10300H
with Python 3.11.2 and PyTorch 2.5.1; both CPUs have 4 cores and 8 threads, both machines
used Gymnasium 1.3.0 and NumPy 2.4.6, and all networks ran on the CPU. All comparisons
within a tier, including wall-clock times, come from a single machine. Tier~1 runs record
commit \texttt{6398ca6} (tag \texttt{v1-experiments}); Tiers~2--3 ran from commit
\texttt{a1c62ff}, which differs only in not recording the commit in the run metadata. One
command per experiment reproduces every figure (see the repository README).

% ---------------------------------------------------------------- 2.3 pages
\section{Results}\label{sec:results}

Table~\ref{tab:results} summarizes every variant. We report two success measures:
\emph{reached}, a 10-episode evaluation during training at or above the threshold, and
\emph{solved}, a final 100-episode mean at or above it, which is the official criterion.

\begin{table}[t]
\centering
\caption{Summary over seeds. Final: return of the final greedy policy over 100 episodes,
mean $\pm$ s.d.\ across seeds. Solved: seeds whose final return reaches the threshold.
Reached: seeds with a 10-episode evaluation at or above it during training; Steps: median
first step at which this happens. Time: mean wall-clock per run.}
\label{tab:results}
\footnotesize
\setlength{\tabcolsep}{2.5pt}
\begin{tabular}{@{}lrccrr@{}}
\toprule
Variant & Final return & Solved & Reached & Steps & Time (s) \\
\midrule
\multicolumn{6}{@{}l}{\emph{Tier 1: CartPole-v1, threshold 475, 5 seeds}} \\
Q-learning, 6 bins       & $216.9\pm53.8$  & 0/5 & 0/5 & --   & 44 \\
Q-learning, 10 bins      & $199.2\pm62.0$  & 0/5 & 0/5 & --   & 50 \\
Q-learning, 20 bins      & $163.6\pm47.9$  & 0/5 & 0/5 & --   & 45 \\
Q-learning, 10 narrow    & $280.2\pm107.1$ & 1/5 & 1/5 & 145k & 46 \\
Q-learning, 10 wide      & $175.5\pm29.0$  & 0/5 & 0/5 & --   & 44 \\
SARSA, 10 bins           & $208.1\pm70.4$  & 0/5 & 0/5 & --   & 44 \\
DQN                      & $159.2\pm68.1$  & 0/5 & 5/5 & 90k  & 823 \\
\midrule
\multicolumn{6}{@{}l}{\emph{Tier 2: LunarLander-v3, threshold 200, 3 seeds}} \\
DQN default              & $206.6\pm30.2$  & 2/3 & 3/3 & 220k & 689 \\
lr $10^{-4}$             & $118.7\pm90.7$  & 0/3 & 2/3 & 230k & 764 \\
lr $2\times10^{-3}$      & $112.9\pm205.1$ & 2/3 & 3/3 & 80k  & 788 \\
$C=250$                  & $152.1\pm88.1$  & 1/3 & 3/3 & 260k & 680 \\
$C=4000$                 & $233.7\pm28.3$  & 2/3 & 3/3 & 250k & 663 \\
No target network        & $82.7\pm76.0$   & 0/3 & 3/3 & 100k & 669 \\
No replay buffer         & $-76.4\pm77.1$  & 0/3 & 2/3 & 265k & 656 \\
\midrule
\multicolumn{6}{@{}l}{\emph{Tier 3: LunarLander-v3, threshold 200, 3 seeds}} \\
Double DQN               & $220.0\pm53.9$  & 2/3 & 3/3 & 180k & 681 \\
Dueling DQN              & $222.7\pm47.2$  & 2/3 & 3/3 & 230k & 738 \\
Double + Dueling         & $214.4\pm49.3$  & 2/3 & 3/3 & 230k & 759 \\
\bottomrule
\end{tabular}
\end{table}

\subsection{Tier 1: tabular RL versus DQN on CartPole-v1}
Fig.~\ref{fig:t1-main} compares Q-learning with 6, 10 and 20 bins with DQN. Per
environment step the tables learn first: with 10 bins the mean evaluation return is 208 at
20k steps, against 122 for DQN. All tabular variants then plateau between about 150 and 280
for the rest of the budget. DQN overtakes them after about 85k steps, scores 500 in at least
one evaluation in every seed, and its mean evaluation return peaks at 426 at 130k steps. It
does not hold this level: the band widens to more than $\pm150$, and the final policies
return only $159\pm68$, below every tabular variant. Per second the order reverses: a
tabular run processes its 200k steps in 44--50\,s (about 4\,400 steps/s), DQN in 823\,s
(about 240 steps/s). The 10-bin table first reaches an evaluation return of 200 after a
median of 6\,s; DQN after 145\,s, and 475 after 340--390\,s. The training-return panel
favours DQN early mainly because its $\varepsilon$ reaches 0.05 at 20k steps, against 100k
for the tables.

Fewer bins gave higher final returns: $217\pm54$ with 6 bins, $199\pm62$ with 10 and
$164\pm48$ with 20; the ordering is monotonic but each gap is smaller than one s.d.
At the end of training the agents had visited 28\% of the 1\,296 states with 6 bins, 11\%
of the $10^4$ with 10 bins and 2.4\% of the $1.6\times10^5$ with 20 bins, i.e., about 370,
1\,060 and 3\,800 states, which received on average 540, 190 and 50 updates each.
Fig.~\ref{fig:t1-sens} varies the velocity ranges at 10 bins. Narrow ranges give the best
tabular result ($280\pm107$), including the only tabular seed that solves the task (final
return 493), and raise the visited fraction to 17\%; wide ranges give $176\pm29$ and lower
it to 5.8\%. SARSA ($208\pm70$) and Q-learning ($199\pm62$) with 10 bins are
indistinguishable, and their curves overlap throughout (figure in the repository).

\begin{figure*}[t]
\centering
\resfig{tier1/figures/tabular_vs_dqn.png}
\caption{Tier~1, CartPole-v1: tabular Q-learning (6, 10, 20 bins) versus DQN under the same
200k-step budget. Left: greedy evaluation return (10 episodes every 5k steps); right:
training return. Mean over 5 seeds $\pm1$ s.d.; dashed line: success threshold (475).}
\label{fig:t1-main}
\end{figure*}

\begin{figure*}[t]
\centering
\resfig{tier1/figures/tabular_ranges.png}
\caption{Tier~1, CartPole-v1: Q-learning with 10 bins and three ranges for the cart
velocity ($v$) and the pole angular velocity ($w$). Mean over 5 seeds $\pm1$ s.d.}
\label{fig:t1-sens}
\end{figure*}

\subsection{Tier 2: DQN on LunarLander-v3}
Fig.~\ref{fig:t2-hp} shows the hyperparameter study. The default DQN reaches 200 in every
seed (median 220k steps) and ends at $207\pm30$, with two of three seeds solving the task.
With a learning rate of $2\times10^{-3}$ learning is fastest (median 80k steps) and two seeds
end at 257 and 259, among the best final returns of Tier~2, but the third collapses after
230k steps: its evaluation return stays between $-420$ and $-900$ from 260k steps on and its
final return is $-177$. Its mean $Q$ (the average of $Q(s,a;\theta)$ over the training
minibatches) climbs from 58 to 252 within 70k steps and then swings down to $-59$, whereas
it never exceeds 58 in the other two seeds. With $10^{-4}$ learning is slower and less reliable: one seed
never exceeds $-9$, and the mean $Q$ decays to about 3 by the end. The target period mainly
affects stability: $C=4000$ gives the highest and most consistent final return
($234\pm28$) and the lowest, smoothest mean $Q$ (at most 36), but reaches the threshold
30k steps later than the default; $C=250$ ends at $152\pm88$, with higher mean $Q$ (up to
64) and one seed that drops from 270 to 33 in its last evaluation.

Fig.~\ref{fig:t2-abl} shows the ablation. Without a target network the agent reaches 200
in every seed, even earlier than full DQN (median 100k steps), but cannot keep it: its
evaluations alternate between about 270 and 100 or less, no seed solves the task, and the
final return is $83\pm76$. Its mean $Q$ is erratic, from $-26$ at 20k steps to peaks of 121
and 155, against at most 56 for full DQN. Without a replay buffer, training returns stay
below about 100, evaluations collapse suddenly (one seed scores $-3\,782$ at 370k steps),
and the final return is $-76\pm77$, the worst of all LunarLander runs.

\begin{figure*}[t]
\centering
\resfig{tier2/figures/learning_rate.png}\\[2pt]
\resfig{tier2/figures/target_update.png}
\caption{Tier~2, LunarLander-v3: DQN sensitivity to the learning rate (top) and to the
target-network update period $C$ (bottom). Left: greedy evaluation return; right: training
return. Mean over 3 seeds $\pm1$ s.d.; dashed line: threshold (200).}
\label{fig:t2-hp}
\end{figure*}

\begin{figure*}[t]
\centering
\resfig{tier2/figures/ablation.png}
\caption{Tier~2 ablation on LunarLander-v3: full DQN, without target network, and without
replay buffer. The large negative excursions are single-seed collapses of the agent without
replay. Mean over 3 seeds $\pm1$ s.d.}
\label{fig:t2-abl}
\end{figure*}

\subsection{Tier 3: Double and Dueling DQN}
In Fig.~\ref{fig:t3} the four curves overlap within their bands. Double ($220\pm54$),
Dueling ($223\pm47$) and Double+Dueling DQN ($214\pm49$) end slightly above vanilla DQN
($207\pm30$), and each solves the task in two of three seeds, as vanilla DQN does. Double
DQN reaches the threshold first (median 180k steps, against 220k). The final mean $Q$ is
32 for Double DQN and 35 for the other three variants.

\begin{figure*}[t]
\centering
\resfig{tier3/figures/dqn_variants.png}
\caption{Tier~3, LunarLander-v3: vanilla, Double, Dueling and Double+Dueling DQN. The
vanilla runs are the Tier~2 default runs. Mean over 3 seeds $\pm1$ s.d.}
\label{fig:t3}
\end{figure*}

% ---------------------------------------------------------------- 1.0 page
\section{Discussion}\label{sec:discussion}

\subsection{Which method learned faster in Tier~1, and why?}
The answer depends on the unit. Per second, the table wins by far: a Q-learning update reads
and writes one entry, whereas DQN runs a forward and backward pass of the network on a
minibatch of 64 at every step, about 18 times the cost per step on our CPU. Per environment
step, the table is faster only at first. It shares nothing between states, but with 6 or
10 bins only 370--1\,060 states are ever visited and each is updated hundreds of times, so
the visited entries settle quickly. The table then stops improving: its plateau around 200
is reached by 20k steps and does not move during the remaining 180k, although
$\varepsilon$ keeps decreasing. This is a limit of the representation, not of the data. A
grid cell merges states whose best action differs, e.g., a pole leaning slightly left that
is moving left or right within the same velocity bin, so the discretized problem is not
Markov and no table over it balances for 500 steps reliably. DQN takes the continuous
observation as input and generalizes across nearby states, so it can represent a finer
decision boundary; fitting it takes more steps, but it is the only method that reaches 475.

DQN's late degradation is consistent with a known weakness of DQN on CartPole: once the
policy balances for 500 steps, the 50k-transition buffer holds about 100 near-identical
successful episodes, transitions close to failure are flushed out, and the network is
trained only on states it already handles, so its values for the states that precede a
fall drift. The logged mean $Q$ is in line with this reading: it rises steadily to 89--96
at 200k steps, close to $1/(1-\gamma)=100$, the value of a pole that never falls (time-limit
cuts bootstrap). Because we evaluate the final policy, this degradation is counted; a
best-checkpoint report would have hidden it.

\subsection{Sensitivity to the key design choice}
For the tables the key choice is the discretization, and two effects compete. More bins
give finer resolution but more entries to fill: with 20 bins the 3\,800 visited states
receive about 50 updates each, against 540 with 6 bins, and with $\alpha=0.1$ an entry needs
tens of updates to forget its initial value and average its noisy targets, as
Section~\ref{sec:pertinence} anticipated from the table size. The ranges matter as much as
the count. With wide ranges most observed velocities fall into a few central bins, so only
5.8\% of the table is used and the effective resolution is coarse; narrow ranges spend the
10 bins where the velocities actually are and clip only rare extreme values. What helps is
resolution where the data lie, not the number of bins. SARSA and Q-learning coincide because
their targets differ only in $Q(s',a')$ versus $\max_{a'}Q(s',a')$: once $\varepsilon=0.05$,
95\% of next actions are greedy, and CartPole has no state where one exploratory action is
so costly that an on-policy learner would prefer a safer path, as in the cliff-walking
example of~\cite{sutton2018}.

For DQN the learning rate trades speed for stability. A step of $2\times10^{-3}$ fits the
targets faster, but it also moves the Q-function far between target syncs; in one seed the
loop through the bootstrap target amplified the values (mean $Q$ up to 252) and the policy
collapsed. A step of $10^{-4}$ was too small to learn within 400k steps in one seed. The
target period acts in the opposite direction: a longer period keeps the regression target
fixed for longer, which gave lower and smoother $Q$ estimates and the most reliable final
returns, at the cost of propagating value information more slowly. With three seeds, only
the divergence at $2\times10^{-3}$ is large relative to the across-seed spread.

\subsection{What did the ablation reveal?}
Both components matter, and in our setting the replay buffer matters more. Without a target
network the bootstrap target $\max_{a'}Q(s',a';\theta)$ changes with every gradient step,
so each update partly chases its own output. Value information propagates immediately, which
explains the early threshold crossings, but the estimates oscillate (mean $Q$ from $-26$ to
155) and so does the policy. Without replay, every minibatch is the same window of the last
128 consecutive transitions, which are strongly correlated and come from one region of the
state space. Gradient steps fit that region and overwrite what was learned elsewhere, and
old transitions are never revisited, so the agent forgets, e.g., how to land while it is
learning to hover; the result is the sudden collapses in Fig.~\ref{fig:t2-abl} and the
lowest final return. These are the two problems that replay and the target network were
introduced to solve~\cite{mnih2015}.

In Tier~3, Double DQN did not measurably reduce the mean $Q$ (32 against 35) and no variant
improved significantly on vanilla DQN. The overestimation that Double DQN removes grows
with the noise of the estimates and with the number of actions in the
$\max$~\cite{vanhasselt2016}; with four actions, a Huber loss, gradient clipping and a
1000-step target period, the bias of vanilla DQN appears to be small here. Dueling helps
most when many actions have similar values~\cite{wang2016}, which LunarLander's four
engines do not obviously provide. Detecting a gain of 10--15 points against a spread of
30--50 would require many more seeds.

\subsection{Limitations}
Tiers~2 and~3 use three seeds, so differences smaller than the across-seed spread are not
conclusive. The hyperparameter study varies one value at a time around a single default
and ignores interactions. Steps to threshold are measured with 10-episode evaluations,
which are noisier than the 100-episode criterion. We report the final policy rather than
the best checkpoint, so late performance drops are counted. Tier~1 and Tiers~2--3 ran on
different machines and library versions, so wall-clock times are comparable only within a
tier. The 10-episode \emph{reached} criterion overstates success: without a target network
every seed reached 200 but none solved the task. The mean $Q$ is averaged over replayed
state--action pairs and was not compared with Monte Carlo returns, so it shows trends but
not the size of the overestimation bias. In Tier~1 the training returns of the two
families are not directly comparable because their $\varepsilon$ schedules differ.

% ---------------------------------------------------------------- 1.0 page
\section{Analysis of Unsuitable Environments}\label{sec:analysis}

\subsection{Continuous-action locomotion}
BipedalWalker-v3 and the MuJoCo~\cite{todorov2012} locomotion tasks have continuous,
multidimensional action spaces (torques in $[-1,1]^d$, $[-0.4,0.4]^{17}$ for Humanoid) and
observations of 11 to 348 dimensions (Table~\ref{tab:locomotion}). DQN cannot use them
directly: its greedy policy needs $\max_a Q(s,a)$, which over a continuous set is itself an
optimization problem to solve at every step. Discretizing each of the $d$ torque dimensions
into $k$ levels gives $k^d$ joint actions, each with its own network output and its own
value estimate. With $k=11$ (a torque resolution of 0.2), HalfCheetah already needs
$1.8\times10^6$ outputs and Ant $2.1\times10^8$; even $k=5$ leaves Humanoid with
$7.6\times10^{11}$. Most of these actions would never be tried, so their values would be
pure extrapolation, and coarse levels cannot represent the smooth, coordinated torques of
a gait. Factorizing the output per action dimension mitigates the explosion but changes the
algorithm~\cite{tavakoli2018}. Tabular methods fail earlier, on the state side: 10 bins
per dimension give $10^{24}$ states for BipedalWalker and $10^{105}$ for Ant, against
typical budgets of about $10^6$ environment steps.

\begin{table}[t]
\centering
\caption{Action- and state-space sizes after uniform discretization, computed by
instantiating each environment (Gymnasium 1.3.0, MuJoCo 3.17.0).}
\label{tab:locomotion}
\footnotesize
\setlength{\tabcolsep}{3.5pt}
\begin{tabular}{@{}lrrrrr@{}}
\toprule
 & & & \multicolumn{2}{c}{$|\mathcal{A}| = k^d$} & 10-bin \\
\cmidrule(lr){4-5}
Environment & Obs. & $d$ & $k=5$ & $k=11$ & $|\mathcal{S}|$ \\
\midrule
BipedalWalker-v3 & 24              & 4  & 625                   & 14\,641               & $10^{24}$ \\
Hopper-v5        & 11              & 3  & 125                   & 1\,331                & $10^{11}$ \\
HalfCheetah-v5   & 17              & 6  & 15\,625               & $1.8\!\times\!10^{6}$ & $10^{17}$ \\
Ant-v5           & 105$^{\dagger}$ & 8  & $3.9\!\times\!10^{5}$ & $2.1\!\times\!10^{8}$ & $10^{105}$ \\
Humanoid-v5      & 348             & 17 & $7.6\!\times\!10^{11}$ & $5.1\!\times\!10^{17}$ & $10^{348}$ \\
\bottomrule
\multicolumn{6}{@{}l}{\scriptsize $^{\dagger}$Includes 78 contact-force terms; 27 without them.}
\end{tabular}
\end{table}

\subsection{Suitable algorithm families}
These tasks call for actor-critic methods whose actor outputs a continuous action directly,
which removes the maximization over actions, while a critic $Q(s,a)$ takes the action as an
input and generalizes across it.
\emph{PPO}~\cite{schulman2017} is on-policy: a Gaussian policy is improved with a clipped
surrogate objective that limits the size of each update. It is robust and parallelizes
well across environments, but it discards data after each update, so it needs more samples.
\emph{SAC}~\cite{haarnoja2018} is off-policy: a stochastic actor maximizes return plus an
entropy bonus that drives exploration, data are reused from a replay buffer, and two
critics limit overestimation; it is among the most sample-efficient methods on MuJoCo.
\emph{TD3}~\cite{fujimoto2018} extends DDPG~\cite{lillicrap2016} with a deterministic actor,
clipped double Q-learning, target-policy smoothing and delayed actor updates, which counter
the same overestimation bias that Double DQN addresses in the discrete case. SAC and TD3
keep DQN's replay buffer and target networks; what changes is the actor that replaces the
$\arg\max$.

\subsection{High-dimensional observations: Atari pixels}
An observation of ALE/Pong-v5~\cite{bellemare2013} is a $210\times160\times3$ array of 8-bit
values, so there are up to $256^{100800}\approx10^{242751}$ distinct observations; the usual
preprocessing to four stacked $84\times84$ grayscale frames still leaves
$256^{28224}\approx10^{67970}$. A table over this space cannot be stored, and, more
fundamentally, it cannot generalize: each entry is learned only from visits to that exact
image, two frames that differ in one pixel are unrelated entries, and exact frames rarely
repeat, so almost every entry would be visited at most once. Coarsening the pixels either
keeps the dimensionality or erases the small objects (ball, paddles) that decide the game.
A single frame is also not Markov, because it shows no velocity; this is why frames are
stacked. DQN addresses both problems with a convolutional network whose shared weights map
visually similar frames to similar values~\cite{mnih2015}, at the cost of millions of
frames and GPU time.

\subsection{Summary: matching the problem to the algorithm family}
Table~\ref{tab:families} condenses the analysis. Each row adds one difficulty that breaks
the method of the row above: a state space too large to enumerate breaks the table, a
continuous action space breaks the $\arg\max$ of DQN, and raw pixels require a network
that extracts features from images before any value can be estimated.

\begin{table}[t]
\centering
\caption{Which family suits which problem structure.}
\label{tab:families}
\footnotesize
\setlength{\tabcolsep}{3pt}
\begin{tabular}{@{}llll@{}}
\toprule
States & Actions & Family & Example \\
\midrule
Small or discretizable & Discrete   & Q-learning, SARSA & CartPole \\
Continuous             & Discrete   & DQN and variants  & LunarLander \\
Continuous             & Continuous & PPO, SAC, TD3     & Ant \\
Pixels                 & Discrete   & DQN with a CNN    & Pong \\
\bottomrule
\end{tabular}
\end{table}

% ---------------------------------------------------------------- 0.2 page
\section{AI-Use Statement}
Generative AI was used under the course policy. A.~Chavez used Claude (Anthropic) through
Claude Code to explain the assignment, review the code for correctness, check the software
environment, compute the environment specifications in Tables~\ref{tab:envs}
and~\ref{tab:locomotion}, extract the statistics in Section~\ref{sec:results} from the
logged result files, and draft the text of this report. J.~Reinoso used Claude (Anthropic)
through Claude Code to set up the repository and the Python environment (including the
Box2D build), to write the agents (\texttt{agents/}, \texttt{envs/}) and the experiment
scripts and configurations (\texttt{experiments/}), to integrate Weights \& Biases logging,
to run smoke tests and timing measurements, and to fill in parts of this report (the
environment choice, this statement and the prompt appendix). Both authors reviewed the
code and must be able to explain every line of it. The authors verified all
AI-generated content: environment dimensions were obtained by instantiating each
environment, every reference was checked against its source~\todo{do it before submission},
and every number in Section~\ref{sec:results} comes from the logged result files. The most
important prompts are listed in Appendix~\ref{app:prompts}.

\bibliographystyle{IEEEtran}
\bibliography{references}

\appendices
\section{Most Important Prompts}\label{app:prompts}
Prompts that shaped the work, translated from Spanish and lightly edited for spelling
(two consecutive messages are merged where one request spanned both), each followed by
how its output was used and checked.

\textbf{A.~Chavez:}
\begin{itemize}
  \item ``Explain everything this assignment says.''
  \item ``Review the files my teammate pushed and check whether my environment works for
        them.''
  \item ``Build the full report skeleton and fill in the parts that do not depend on
        results.''
\end{itemize}

\textbf{J.~Reinoso:}
\begin{enumerate}
  \item \textit{``Explain Tier~2 to me. Do you think it can run on my laptop?''}\\
        \textbf{Use and verification:} the answer only estimated the cost. We then timed
        30k LunarLander steps with four runs in parallel on the i5-10210U (about 73\,s per
        run), extrapolated about 20 minutes per 400k-step run, and kept LunarLander instead
        of Acrobot (Section~\ref{sec:pertinence}).
  \item \textit{``Aldrin gave me an idea so we do not step on each other's code: one of us
        takes care of the whole agents part, the tabular agent and the DQN for both tiers,
        and the other one launches the runs. What do you think, or what issues do you see?
        I liked it, but we are leaving out \texttt{common} and \texttt{plot}.''}\\
        \textbf{Use and verification:} we adopted the split (agents and discretizer on one
        side, experiment scripts on the other) and fixed a common agent interface
        (\texttt{act}, \texttt{greedy}, \texttt{update}). Evaluation receives a policy
        function rather than an agent, so the same code evaluates the Q-table and the
        network, which keeps the Tier~1 comparison fair.
  \item \textit{``Generate \texttt{tabular.py} (Q-learning and SARSA),
        \texttt{discretizer.py} and \texttt{dqn.py} with the flags \texttt{use\_target},
        \texttt{use\_replay}, \texttt{double} and \texttt{dueling}. Keep in mind that we
        want it to run in parallel and use the cores; both PCs have 8 cores.''}\\
        \textbf{Use and verification:} the agents were run with all 16 flag combinations
        and a save/load round trip. A parallel test showed that the ``8 cores'' are 4
        physical cores with 2 threads each: 8 parallel CartPole runs took 284\,s against
        75\,s for one run alone, only about $2\times$ faster than running them one after
        another, so every sweep uses 4 processes with one PyTorch thread each.
  \item \textit{``I am using CoreWeave [Forge, a Weights \& Biases server]; I would like
        to upload the plots there.''}\\
        \textbf{Use and verification:} every run writes CSV logs, which are the source of
        all figures, and mirrors them to Weights \& Biases. Uploads were checked with test
        runs in a separate project, listed through the W\&B API, before the real sweeps.
  \item \textit{``Aldrin says we should freeze the commits; it is for W\&B. For that we
        need a line of code in each experiment, adjust it before running.''}\\
        \textbf{Use and verification:} instead of one line per experiment, the logger
        records the commit, tag and uncommitted-change flag of every run, and the code was
        tagged \texttt{v1-experiments}. All 35 Tier~1 summaries report this tag with no
        uncommitted changes; the Tier~2--3 runs predate the feature, and their commit
        (\texttt{a1c62ff}) was recovered from the W\&B metadata.
\end{enumerate}

\end{document}
